\documentclass[12pt, a4paper]{article}

% --- PAQUETES Y CONFIGURACIÓN ---
\usepackage[utf8]{inputenc}
\usepackage[spanish, es-tabla]{babel}
\usepackage[margin=2.5cm, headheight=28pt]{geometry}
\usepackage{graphicx}
\usepackage{xcolor}
\usepackage{fancyhdr}
\usepackage{titlesec}
\usepackage{hyperref}
\usepackage{enumitem}
\usepackage{longtable}
\usepackage{booktabs}
\usepackage{array}

% Color primario institucional
\definecolor{uachpurple}{RGB}{88, 38, 92}
\definecolor{uachdark}{RGB}{40, 40, 40}

% Configuración de hipervínculos
\hypersetup{
    colorlinks=true,
    linkcolor=uachpurple,
    urlcolor=uachpurple,
    citecolor=uachpurple
}

% Encabezados y pies de página
\pagestyle{fancy}
\fancyhf{}
\lhead{\small\color{uachdark} Desarrollo de Aplicaciones Web}
\rhead{\small\color{uachdark} Tarea 2. Sprint 0}
\cfoot{\thepage}
\renewcommand{\headrulewidth}{0.4pt}
\renewcommand{\footrulewidth}{0.4pt}

% Formato de Títulos
\titleformat{\section}
  {\normalfont\Large\bfseries\color{uachpurple}}{\thesection.}{0.5em}{}
\titleformat{\subsection}
  {\normalfont\large\bfseries\color{uachpurple}}{\thesubsection.}{0.5em}{}
\titleformat{\subsubsection}
  {\normalfont\normalsize\bfseries\color{uachpurple}}{\thesubsubsection.}{0.5em}{}

% Atajo para el encabezado de cada Historia de Usuario
\newcommand{\hutag}[3]{\textit{#1 · #2 · #3}\par\vspace{2pt}}

\begin{document}

% --- PORTADA ---
\begin{titlepage}
    \centering
    % Descomentar si deseas incluir el logo oficial de la facultad:
    % \includegraphics[width=0.25\textwidth]{logo_uach.png}\\[1cm]

    {\Large\bfseries Universidad Autónoma de Chihuahua}\\[0.3cm]
    {\large Facultad de Ingeniería}\\[0.2cm]
    {\large Ingeniería en Computación}\\[1.5cm]

    \rule{\linewidth}{1.5pt}\\[0.4cm]
    {\huge\bfseries Constitución del Proyecto}\\[0.2cm]
    {\LARGE (Sprint 0) --- Stride \& Co.}\\[0.4cm]
    \rule{\linewidth}{1.5pt}\\[1.5cm]

    \begin{minipage}{0.48\textwidth}
        \begin{flushleft} \large
            \textbf{Curso:} Desarrollo de Aplicaciones Web\\
            \textbf{Docente:} Mtro. Luis Antonio Ramírez Martínez\\
            \textbf{Fecha:} 28 de agosto de 2026\\
            \textbf{Equipo:} CodeMen
        \end{flushleft}
    \end{minipage}
    \hfill
    \begin{minipage}{0.48\textwidth}
        \begin{flushright} \large
            \textbf{Equipo Scrum:}\\
            Osvaldo Hernández Juárez (\textit{Product Owner})\\
            José Alejandro Pérez Millán (\textit{Scrum Master})\\
            Jenny Guadalupe Quintana Sánchez (\textit{Developer})\\
            Armando Velázquez Barroso (\textit{Developer})
        \end{flushright}
    \end{minipage}

    \vfill
    {\small Chihuahua, Chih., México}
\end{titlepage}

% --- ÍNDICE ---
\newpage
\tableofcontents
\newpage

% --- 1. VISIÓN DEL PRODUCTO ---
\section{Visión del Producto}

\subsection{Problema a Resolver}
Stride \& Co. es una tienda deportiva que gestiona ventas, inventario y pedidos de forma manual, usando WhatsApp y hojas de cálculo. Esto provoca pedidos perdidos entre chats, falta de visibilidad del inventario en tiempo real, ausencia de registros formales de compra y pago, y una sobrecarga de trabajo para el equipo de operaciones y ventas al no poder consultar la información de forma centralizada.

\subsection{Objetivos del Producto}
\begin{itemize}[leftmargin=*]
    \item Digitalizar el flujo completo de ventas: publicación de catálogo, compra, confirmación y seguimiento de pedido.
    \item Automatizar el descuento de inventario al confirmarse un pedido, eliminando el conteo manual.
    \item Brindar visibilidad en tiempo real del stock y del estado de los pedidos a cada rol involucrado.
    \item Establecer roles y permisos claros que segmenten el acceso a la información según el tipo de usuario.
    \item Garantizar que el flujo de compra funcione de manera confiable de principio a fin.
\end{itemize}

\subsection{Usuarios Objetivo}
\begin{itemize}[leftmargin=*]
    \item Clientes finales que desean comprar productos deportivos y dar seguimiento a su pedido sin depender de mensajes directos.
    \item Equipo de Ventas que necesita gestionar los pedidos que le corresponden.
    \item Equipo de Operaciones que necesita visibilidad de inventario y notificación inmediata de nuevos pedidos.
    \item Administración que necesita supervisión completa de la operación.
\end{itemize}

\subsection{Valor de Negocio}
Reducir pedidos perdidos o que no tengan seguimiento, para disminuir las quejas y reclamos de los clientes. Disminuir el tiempo operativo dedicado a responder las constantes preguntas de los clientes respecto a su pedido. Dar trazabilidad formal a cada venta, evitando la pérdida de información. Permite escalar el negocio sin que el crecimiento en ventas implique el caos operativo proporcional.

\subsection{Métricas de Éxito}
\begin{itemize}[leftmargin=*]
    \item \% de pedidos con registro formal completo (meta: 100\%, vs. registros informales actuales).
    \item Tiempo promedio entre la creación de un pedido y su notificación al equipo de operaciones (meta: near-instantáneo).
    \item Reducción de consultas manuales de stock a bodega.
    \item Reducción de mensajes de clientes preguntando por el estado de su pedido.
\end{itemize}


% --- 2. CONSTITUCIÓN DEL EQUIPO SCRUM ---
\section{Constitución del Equipo Scrum}

\subsection{Integrantes}
\begin{center}
\begin{tabular}{ll}
\toprule
\textbf{Integrante} & \textbf{Matrícula} \\
\midrule
Osvaldo Hernández Juárez & 385493 \\
José Alejandro Pérez Millán & 385570 \\
Armando Velázquez Barroso & 254106 \\
Jenny Guadalupe Quintana Sánchez & 374366 \\
\bottomrule
\end{tabular}
\end{center}

\subsection{Roles}
\begin{itemize}[leftmargin=*]
    \item \textbf{Osvaldo Hernández Juárez --- Product Owner:} representa la voz del cliente (Camila, Jorge, Valeria), prioriza y refina el Product Backlog, y decide qué se construye en cada sprint.
    \item \textbf{José Alejandro Pérez Millán --- Scrum Master:} facilita las ceremonias del equipo (planeación, daily, revisión, retrospectiva), da seguimiento a impedimentos y protege el proceso Scrum.
    \item \textbf{Armando Velázquez Barroso --- Developer:} implementa las historias de usuario del backlog y participa en la revisión de código.
    \item \textbf{Jenny Guadalupe Quintana Sánchez --- Developer:} implementa las historias de usuario del backlog y participa en la revisión de código.
\end{itemize}

\subsection{Herramientas}
\begin{itemize}[leftmargin=*]
    \item \textbf{Gestión del proyecto:} GitHub Projects
    \item \textbf{Control de versiones:} Git + GitHub
    \item \textbf{Comunicación:} WhatsApp
    \item \textbf{Documentación:} LaTeX
    \item \textbf{Stack de desarrollo:} MEVN --- MongoDB, Express, Vue, Node.js, con Socket.io para tiempo real y JWT para autenticación y roles.
\end{itemize}

\subsection{Estrategia de Colaboración}
\begin{itemize}[leftmargin=*]
    \item División de trabajo por historias de usuario.
    \item Rama \texttt{main} protegida; desarrollo en ramas por feature.
    \item Pull requests obligatorios antes de hacer merge a \texttt{main}; deben ser revisados por el equipo.
    \item El Scrum Master da seguimiento a los impedimentos detectados y facilita la planeación y revisión de cada sprint.
    \item El Product Owner revisa y prioriza el backlog antes del inicio de cada sprint.
\end{itemize}

\subsection{Organización del Proyecto}
\begin{itemize}[leftmargin=*]
    \item Las historias de usuario del Product Backlog se convierten en Issues vinculados al GitHub Project, utilizando la plantilla \textit{Iterative Development}.
    \item Estructura del repositorio separada por capa (frontend / backend), a definir con más detalle en la primera reunión técnica del equipo.
    \item Ambiente de desarrollo local individual por integrante, con variables de entorno compartidas mediante un archivo \texttt{.env.example}.
    \item Convenciones de commits y de nombres de rama documentadas en el \texttt{README.md} del repositorio para mantener consistencia entre integrantes.
\end{itemize}


% --- 3. ANÁLISIS DEL PROYECTO ---
\section{Análisis del Proyecto}

\subsection{Necesidades}
Digitalizar la venta y eliminar la dependencia de WhatsApp como canal transaccional. Visibilidad de inventario en tiempo real para operaciones. Registro formal y centralizado de cada pedido y su estado de pago. Separación clara de accesos por rol.

\subsection{Actores}
\begin{itemize}[leftmargin=*]
    \item \textbf{Cliente:} navega catálogo, compra, da seguimiento a su pedido.
    \item \textbf{Vendedor:} gestiona únicamente los pedidos que le corresponden.
    \item \textbf{Operaciones:} consulta stock en tiempo real, recibe notificaciones de pedidos nuevos.
    \item \textbf{Administrador:} supervisa toda la operación (pedidos y vendedores).
\end{itemize}

\subsection{Restricciones}
El descuento de inventario debe ser automático al confirmarse un pedido (no negociable). Los roles deben estar estrictamente segmentados: un cliente no accede al catálogo administrativo; un vendedor no ve pedidos de otro vendedor. Las notificaciones y actualizaciones de estado deben ser en tiempo real, no por refresco manual. El proyecto debe desarrollarse en el marco de un semestre académico, con entregas incrementales por sprint.

\subsection{Supuestos}
El negocio ya cuenta con un catálogo de productos (con variantes de talla y modelo) que puede digitalizarse. Existen varios vendedores operando de forma simultánea. Los usuarios finales tienen acceso a internet y un dispositivo con navegador. La primera versión no requiere integración con pasarela de pago real (puede simularse la confirmación de pago).

\subsection{Funcionalidades Esperadas}
\begin{itemize}[leftmargin=*]
    \item Catálogo de productos navegable por el cliente.
    \item Carrito de compra y confirmación de pedido.
    \item Seguimiento de estado de pedido en tiempo real para el cliente.
    \item Panel de vendedor con sus pedidos asignados.
    \item Panel de operaciones con inventario en tiempo real y notificaciones de pedidos nuevos.
    \item Panel administrativo con visibilidad total.
    \item Descuento automático de inventario al confirmar pedido.
    \item Sistema de roles y autenticación.
\end{itemize}


% --- 4. PRODUCT BACKLOG ---
\section{Product Backlog}
\textit{Proyecto Integrador · Desarrollo de Aplicaciones Web · UACH · Equipo CodeMen · Versión inicial (Sprint 0)}

\subsection{Épicas}
\begin{longtable}{@{}p{1.3cm}p{4.3cm}p{9.5cm}@{}}
\toprule
\textbf{ID} & \textbf{Épica} & \textbf{Objetivo} \\
\midrule
\endhead
E1 & Catálogo y compra autogestionada & Que el cliente compre y dé seguimiento sin escribirle a nadie \\
E2 & Panel de gestión de ventas & Centralizar la trazabilidad de pedidos con visibilidad por rol \\
E3 & Control automatizado de stock & Garantizar la integridad del inventario y eliminar la venta doble \\
E4 & Gestión de roles y seguridad & Jerarquía estricta de permisos y aislamiento de datos \\
E5 & Notificaciones en tiempo real & Eliminar la espera pasiva del staff y del cliente \\
E6 & Gestión de catálogo & Permitir que la operación publique y mantenga los productos \\
\bottomrule
\end{longtable}

\subsection{Historias de Usuario}

\subsubsection{E1 --- Catálogo y compra autogestionada}

\paragraph{HU-19. Registro e inicio de sesión de cliente}
\hutag{Must}{3 SP}{E1}
Como \textbf{cliente}, quiero crear una cuenta e iniciar sesión, para dar seguimiento a mis pedidos y recibir avisos.
\begin{itemize}[leftmargin=*]
    \item \textbf{Dado} que soy un visitante nuevo, \textbf{cuando} me registro con correo y contraseña, \textbf{entonces} el sistema crea mi cuenta con rol Cliente y la contraseña queda almacenada cifrada.
\end{itemize}

\paragraph{HU-01. Visualización del catálogo}
\hutag{Must}{5 SP}{E1}
Como \textbf{cliente}, quiero navegar el catálogo con modelos y tallas, para conocer lo disponible sin preguntar por chat.
\begin{itemize}[leftmargin=*]
    \item \textbf{Dado} que accedo a la plataforma, \textbf{cuando} abro el catálogo, \textbf{entonces} veo los productos publicados con sus tallas y la disponibilidad real tomada de la base de datos.
    \item \textbf{Dado} que una talla tiene stock cero, \textbf{cuando} la visualizo, \textbf{entonces} aparece marcada como agotada y no puedo seleccionarla.
\end{itemize}

\paragraph{HU-02. Carrito de compra}
\hutag{Must}{5 SP}{E1}
Como \textbf{cliente}, quiero agregar productos a un carrito, para comprar varios artículos en una sola operación.
\begin{itemize}[leftmargin=*]
    \item \textbf{Dado} que selecciono un producto y talla, \textbf{cuando} lo agrego al carrito, \textbf{entonces} se registra con la cantidad elegida y puedo modificarla o eliminarlo.
    \item \textbf{Dado} que solicito más unidades de las disponibles, \textbf{cuando} intento agregarlas, \textbf{entonces} el sistema me lo impide indicando el máximo.
\end{itemize}

\paragraph{HU-03. Confirmación de pedido}
\hutag{Must}{5 SP}{E1}
Como \textbf{cliente}, quiero confirmar mi compra, para que quede un pedido formal y deje de depender de una imagen en el chat.
\begin{itemize}[leftmargin=*]
    \item \textbf{Dado} que tengo productos en el carrito, \textbf{cuando} confirmo, \textbf{entonces} el sistema genera un pedido con identificador único, fecha, mis datos, el detalle de artículos y el importe total, en estado \textit{Pendiente de pago}.
    \item \textbf{Dado} que el pedido se creó, \textbf{cuando} consulto la base de datos, \textbf{entonces} existe un registro persistente e íntegro de la venta.
\end{itemize}

\paragraph{HU-04. Seguimiento del pedido}
\hutag{Must}{5 SP}{E1}
Como \textbf{cliente}, quiero consultar el estado de mi pedido, para no tener que preguntar si ya casi me llega.
\begin{itemize}[leftmargin=*]
    \item \textbf{Dado} que inicié sesión, \textbf{cuando} entro a mis pedidos, \textbf{entonces} veo cada uno con su identificador, estado actual y fecha de la última actualización.
\end{itemize}

\subsubsection{E6 --- Gestión de catálogo}

\paragraph{HU-14. Publicación de producto}
\hutag{Must}{5 SP}{E6}
Como \textbf{administrador}, quiero publicar un producto nuevo con sus tallas, precio y stock inicial, para que aparezca en el catálogo.
\begin{itemize}[leftmargin=*]
    \item \textbf{Dado} que tengo rol Administrador, \textbf{cuando} registro un producto con al menos una talla y su stock, \textbf{entonces} queda publicado y visible para los clientes.
    \item \textbf{Dado} que dejo campos obligatorios vacíos, \textbf{cuando} intento guardar, \textbf{entonces} el sistema rechaza el alta y señala qué falta.
\end{itemize}

\paragraph{HU-15. Edición y despublicación}
\hutag{Must}{3 SP}{E6}
Como \textbf{administrador}, quiero editar o despublicar un producto, para mantener el catálogo actualizado sin perder el histórico de ventas.
\begin{itemize}[leftmargin=*]
    \item \textbf{Dado} que un producto tiene pedidos asociados, \textbf{cuando} lo despublico, \textbf{entonces} deja de aparecer en el catálogo pero sus pedidos previos conservan la información del artículo.
\end{itemize}

\subsubsection{E3 --- Control automatizado de stock}

\paragraph{HU-08. Consulta de stock en tiempo real}
\hutag{Must}{5 SP}{E3}
Como \textbf{responsable de operaciones}, quiero consultar el stock por modelo y talla, para dejar de llamar físicamente a la bodega.
\begin{itemize}[leftmargin=*]
    \item \textbf{Dado} que inicié sesión con rol de operaciones, \textbf{cuando} abro el inventario, \textbf{entonces} veo las existencias actuales por modelo y talla, diferenciando unidades disponibles de reservadas.
\end{itemize}

\paragraph{HU-09. Reserva y descuento automático de stock}
\hutag{Must (no negociable)}{8 SP}{E3}
Como \textbf{responsable de operaciones}, quiero que el stock se aparte al crear el pedido y se descuente al confirmarse el pago, para que nunca se venda dos veces el mismo artículo.
\begin{itemize}[leftmargin=*]
    \item \textbf{Dado} que un cliente confirma un pedido, \textbf{cuando} el sistema lo registra, \textbf{entonces} las unidades quedan reservadas y dejan de ofrecerse a otros clientes.
    \item \textbf{Dado} que dos clientes intentan comprar la última unidad simultáneamente, \textbf{cuando} ambos confirman, \textbf{entonces} solo uno obtiene la reserva y el otro recibe aviso de falta de disponibilidad.
    \item \textbf{Dado} que un vendedor marca el pedido como pagado, \textbf{cuando} se guarda el cambio, \textbf{entonces} la reserva se convierte en descuento definitivo del inventario.
    \item \textbf{Dado} que una reserva supera el plazo definido sin pago, \textbf{cuando} expira, \textbf{entonces} las unidades regresan al stock disponible.
\end{itemize}

\paragraph{HU-16. Registro de entrada de mercancía}
\hutag{Must}{3 SP}{E3}
Como \textbf{responsable de operaciones}, quiero registrar entradas y ajustes de inventario, para que el stock del sistema refleje lo que hay en bodega.
\begin{itemize}[leftmargin=*]
    \item \textbf{Dado} que llega mercancía nueva, \textbf{cuando} registro las unidades por talla, \textbf{entonces} el stock disponible aumenta y queda constancia de quién hizo el ajuste y cuándo.
\end{itemize}

\paragraph{HU-18. Cancelación con devolución de stock}
\hutag{Should}{3 SP}{E3}
Como \textbf{vendedor}, quiero cancelar un pedido y que las unidades vuelvan al inventario, para que el stock no se degrade con pedidos que no se concretaron.
\begin{itemize}[leftmargin=*]
    \item \textbf{Dado} que un pedido está reservado o pagado, \textbf{cuando} lo cancelo, \textbf{entonces} las unidades regresan al stock disponible y el pedido queda en estado \textit{Cancelado}.
\end{itemize}

\subsubsection{E4 --- Gestión de roles y seguridad}

\paragraph{HU-10. Autenticación con roles diferenciados}
\hutag{Must}{5 SP}{E4}
Como \textbf{administrador}, quiero que cada usuario acceda según su rol, para proteger la operación del negocio.
\begin{itemize}[leftmargin=*]
    \item \textbf{Dado} que un usuario inicia sesión, \textbf{cuando} se valida su identidad, \textbf{entonces} accede únicamente a los módulos correspondientes a su rol: Cliente, Vendedor o Administrador.
    \item \textbf{Dado} que soy Cliente, \textbf{cuando} intento acceder a la gestión del catálogo, \textbf{entonces} el sistema deniega la petición.
\end{itemize}

\paragraph{HU-11. Aislamiento de pedidos por vendedor}
\hutag{Must}{5 SP}{E4}
Como \textbf{vendedor}, quiero que solo yo vea y modifique mis pedidos, para que nadie más altere mi trabajo.
\begin{itemize}[leftmargin=*]
    \item \textbf{Dado} que estoy autenticado como Vendedor A, \textbf{cuando} consulto mi panel, \textbf{entonces} veo exclusivamente los pedidos asignados a mi cuenta.
    \item \textbf{Dado} que soy Vendedor A, \textbf{cuando} accedo por URL directa a un pedido del Vendedor B, \textbf{entonces} el sistema responde con acceso denegado (403).
    \item \textbf{Dado} que soy Administrador, \textbf{cuando} consulto pedidos, \textbf{entonces} veo los de todos los vendedores.
\end{itemize}

\subsubsection{E2 --- Panel de gestión de ventas}

\paragraph{HU-17. Asignación automática rotativa de pedidos}
\hutag{Must}{3 SP}{E2}
Como \textbf{administrador}, quiero que los pedidos web se repartan automáticamente entre los vendedores activos, para que cada pedido tenga un responsable claro.
\begin{itemize}[leftmargin=*]
    \item \textbf{Dado} que entra un pedido desde la web, \textbf{cuando} el sistema lo registra, \textbf{entonces} lo asigna por turno rotativo a un vendedor activo y ese vendedor queda como responsable.
\end{itemize}

\paragraph{HU-05. Panel de pedidos del vendedor}
\hutag{Must}{8 SP}{E2}
Como \textbf{vendedor}, quiero un panel con mis pedidos clasificados por estado, para gestionar mi trabajo de forma ordenada.
\begin{itemize}[leftmargin=*]
    \item \textbf{Dado} que inicié sesión como Vendedor, \textbf{cuando} abro mi panel, \textbf{entonces} veo mis pedidos con identificador, cliente, importe y estado, y puedo filtrarlos entre pagados y pendientes.
\end{itemize}

\paragraph{HU-06. Actualización del estado de un pedido}
\hutag{Must}{3 SP}{E2}
Como \textbf{vendedor}, quiero cambiar el estado de mis pedidos, para mantener informados al equipo y al cliente.
\begin{itemize}[leftmargin=*]
    \item \textbf{Dado} que gestiono un pedido, \textbf{cuando} cambio su estado, \textbf{entonces} el sistema guarda el nuevo estado junto con el usuario responsable y la fecha del cambio, y dispara las notificaciones correspondientes.
\end{itemize}

\paragraph{HU-07. Dashboard global de administración}
\hutag{Should}{5 SP}{E2}
Como \textbf{dueña del negocio}, quiero una vista global de ventas y stock, para tomar decisiones sin preguntarle a cada vendedor.
\begin{itemize}[leftmargin=*]
    \item \textbf{Dado} que inicié sesión como Administrador, \textbf{cuando} abro el dashboard, \textbf{entonces} veo el total de pedidos por estado, las ventas agregadas de todos los vendedores y los productos con inventario crítico.
\end{itemize}

\subsubsection{E5 --- Notificaciones en tiempo real}

\paragraph{HU-12. Alerta de pedido nuevo}
\hutag{Should}{5 SP}{E5}
Como \textbf{responsable de operaciones}, quiero recibir un aviso instantáneo al entrar un pedido, para no estar refrescando la pantalla cada cinco minutos.
\begin{itemize}[leftmargin=*]
    \item \textbf{Dado} que tengo la plataforma abierta, \textbf{cuando} un cliente confirma un pedido, \textbf{entonces} aparece una alerta visible en menos de cinco segundos sin que yo recargue la página.
\end{itemize}

\paragraph{HU-13. Notificación de cambio de estado al cliente}
\hutag{Should}{5 SP}{E5}
Como \textbf{cliente}, quiero enterarme cuando mi pedido cambia de estado, para no tener que preguntar.
\begin{itemize}[leftmargin=*]
    \item \textbf{Dado} que tengo un pedido activo, \textbf{cuando} el vendedor cambia su estado, \textbf{entonces} recibo la actualización en mi cuenta sin recargar la página.
\end{itemize}

\subsection{Priorización Inicial}
El criterio de orden es habilitar primero lo que sostiene al resto del sistema (autenticación y catálogo), continuar con el flujo de venta y la regla no negociable del control de stock, y dejar al final las funcionalidades que enriquecen la experiencia. Esta secuencia coincide con la organización del semestre, que ubica la lógica de servidor y la persistencia de datos en los Sprints 2 a 4, y la comunicación en tiempo real en los Sprints 5 y 6.

\begin{longtable}{@{}p{1cm}p{1.3cm}p{5.2cm}p{2cm}p{1cm}p{1.3cm}@{}}
\toprule
\textbf{Orden} & \textbf{ID} & \textbf{Historia} & \textbf{MoSCoW} & \textbf{SP} & \textbf{Sprint} \\
\midrule
\endhead
1 & HU-10 & Autenticación con roles & Must & 5 & 1 \\
2 & HU-19 & Registro de cliente & Must & 3 & 1 \\
3 & HU-14 & Publicar producto & Must & 5 & 1 \\
4 & HU-01 & Ver catálogo & Must & 5 & 1 \\
5 & HU-15 & Editar o despublicar producto & Must & 3 & 2 \\
6 & HU-08 & Consultar stock & Must & 5 & 2 \\
7 & HU-02 & Carrito de compra & Must & 5 & 2 \\
8 & HU-03 & Confirmar pedido & Must & 5 & 2 \\
9 & HU-09 & Reserva y descuento de stock & Must & 8 & 3 \\
10 & HU-17 & Asignación rotativa & Must & 3 & 3 \\
11 & HU-05 & Panel del vendedor & Must & 8 & 3 \\
12 & HU-06 & Actualizar estado de pedido & Must & 3 & 3 \\
13 & HU-11 & Aislamiento por vendedor & Must & 5 & 4 \\
14 & HU-04 & Seguimiento del pedido & Must & 5 & 4 \\
15 & HU-16 & Entrada de mercancía & Must & 3 & 4 \\
16 & HU-18 & Cancelar y devolver stock & Should & 3 & 4 \\
17 & HU-12 & Alerta de pedido nuevo & Should & 5 & 5 \\
18 & HU-13 & Notificación al cliente & Should & 5 & 5 \\
19 & HU-07 & Dashboard global & Should & 5 & 6 \\
\bottomrule
\end{longtable}

\textbf{Total: 89 story points} --- 13 historias \textit{Must have} y 6 historias \textit{Should have}.

\paragraph{Fuera de alcance (Won't have)}
Consideradas y descartadas conscientemente en esta versión:
\begin{itemize}[leftmargin=*]
    \item Pasarela de pagos en línea
    \item Notificaciones por correo electrónico o WhatsApp
    \item Número de guía de envío
    \item Gestión de devoluciones
\end{itemize}

\subsection{Requerimientos No Funcionales}
\begin{longtable}{@{}p{1.4cm}p{2.3cm}p{10.5cm}@{}}
\toprule
\textbf{ID} & \textbf{Atributo} & \textbf{Especificación} \\
\midrule
\endhead
RNF-01 & Integridad de datos & El descuento de inventario debe ejecutarse de forma atómica: dos pedidos simultáneos sobre la última unidad nunca pueden concretarse ambos. \\
RNF-02 & Seguridad & Control de acceso basado en roles con jerarquía Cliente $<$ Vendedor $<$ Administrador, aplicado a nivel de API y no solo de interfaz. \\
RNF-03 & Seguridad & Contraseñas almacenadas mediante hash y comunicación cifrada con HTTPS de extremo a extremo. \\
RNF-04 & Rendimiento & La actualización de stock debe reflejarse en menos de 2 segundos para evitar condiciones de carrera. \\
RNF-05 & Rendimiento & Las notificaciones llegan en menos de 5 segundos sin recarga manual; se implementan con WebSockets y no con sondeo periódico. \\
RNF-06 & Confiabilidad & Las pruebas se realizan con datos representativos del catálogo e historial real, no con datos ficticios de demostración. \\
RNF-07 & Estabilidad & El flujo completo (publicar $\rightarrow$ comprar $\rightarrow$ actualizar estado $\rightarrow$ notificar) debe ejecutarse sin fallos bajo carga concurrente equivalente a la operación actual. \\
RNF-08 & Trazabilidad & Todo cambio de estado de un pedido registra usuario responsable, fecha y estado anterior. \\
RNF-09 & Usabilidad & La plataforma debe ser utilizable desde dispositivo móvil, dado que el canal actual del cliente es el teléfono. \\
RNF-10 & Mantenibilidad & Código versionado en Git con historial de commits descriptivo y separación clara entre capas. \\
\bottomrule
\end{longtable}

Los requerimientos RNF-01, RNF-04, RNF-05, RNF-06 y RNF-07 se derivan directamente de la entrevista con el cliente; los restantes corresponden a estándares asumidos por el equipo.

\subsection{Riesgos}
\begin{longtable}{@{}p{1.2cm}p{4.5cm}p{1.5cm}p{1.7cm}p{5.5cm}@{}}
\toprule
\textbf{ID} & \textbf{Riesgo} & \textbf{Prob.} & \textbf{Impacto} & \textbf{Mitigación} \\
\midrule
\endhead
R-01 & Condición de carrera en el descuento de stock que permite la venta doble & Alta & Crítico & Reserva atómica al confirmar el pedido y pruebas explícitas de concurrencia sobre HU-09 \\
R-02 & La curva de aprendizaje del stack MEVN retrasa los primeros sprints & Alta & Alto & Sprint 1 dedicado a fundamentos y programación en pareja entre developers \\
R-03 & MongoDB no ofrece transacciones multi-documento de forma natural & Media & Alto & Modelar el stock dentro del documento de producto y usar operaciones atómicas condicionales \\
R-04 & Los WebSockets se subestiman y las notificaciones no llegan a tiempo & Media & Medio & Programarlos en el Sprint 5 con margen y mantener el sondeo periódico como plan alterno \\
R-05 & Crecimiento no controlado del alcance durante el semestre & Media & Alto & Backlog congelado por sprint; los cambios ingresan al backlog y no al sprint en curso \\
R-06 & La carga académica del equipo reduce la capacidad real por sprint & Alta & Medio & Estimación conservadora y monitoreo de la velocidad real por el Scrum Master desde el Sprint 1 \\
R-07 & Ambigüedad en requisitos que el cliente no especificó & Media & Medio & Documentar cada supuesto y validarlo durante la Sprint Review \\
R-08 & No se cuenta con datos reales del catálogo pese a que el cliente los exige & Media & Medio & Construir un conjunto de datos representativo y solicitar el histórico con anticipación \\
R-09 & Dependencia de un solo integrante en un área del sistema & Media & Alto & Rotación de responsabilidades y revisión cruzada de código \\
\bottomrule
\end{longtable}

\subsection{Deuda Técnica Inicial}
El equipo asume de forma deliberada la siguiente deuda técnica al inicio del proyecto, con el compromiso de resolverla en los sprints indicados.

\begin{longtable}{@{}p{1.4cm}p{6.8cm}p{6.2cm}@{}}
\toprule
\textbf{ID} & \textbf{Deuda técnica} & \textbf{Plan de resolución} \\
\midrule
\endhead
DT-01 & El equipo no domina aún Vue ni Mongoose, por lo que el código inicial será mejorable & Refactorización tras el Sprint 2, una vez estabilizado el aprendizaje \\
DT-02 & El esquema de datos no está definido y se modelará sobre la marcha & Formalizar el modelo al cerrar el Sprint 1 \\
DT-03 & Sin pipeline de integración continua ni despliegue automatizado & Incorporarlo a partir del Sprint 3 \\
DT-04 & Cobertura de pruebas automatizadas nula al inicio & Priorizar pruebas sobre HU-09 y HU-11, por ser las críticas \\
DT-05 & Sin ambiente productivo definido; únicamente desarrollo local & Definir el hosting antes del Sprint 5 \\
DT-06 & Imágenes de producto manejadas por URL externa, sin gestión de archivos & Evaluar almacenamiento propio si el cliente lo requiere \\
DT-07 & Sin sistema de registro de eventos ni monitoreo & Integrar logging al trabajar la trazabilidad (RNF-08) \\
DT-08 & Interfaz sin sistema de diseño; se emplearán componentes base & Unificar estilos en el Sprint 5, junto con la experiencia de cliente \\
\bottomrule
\end{longtable}


% --- 5. REPOSITORIO DEL PROYECTO (GITHUB) ---
\section{Repositorio del Proyecto (GitHub)}

\subsection{URL del Repositorio}
% TODO: agregar el enlace real una vez creado el repositorio
\url{https://github.com/usuario/stride-co}

\subsection{Evidencia del Repositorio (README.md)}
% TODO: insertar captura de pantalla del README.md visible en GitHub
% \includegraphics[width=\linewidth]{evidencia_readme.png}

\subsection{Evidencia del GitHub Project}
% TODO: insertar captura de pantalla del GitHub Project (plantilla Iterative Development) mostrando el Product Backlog inicial
% \includegraphics[width=\linewidth]{evidencia_project.png}

\end{document}